\documentclass[12pt,a4paper]{report}

\usepackage[utf8]{inputenc}
\usepackage[T1]{fontenc}
\usepackage[french]{babel}
\usepackage[top=2.5cm,bottom=2.5cm,left=3cm,right=2.5cm]{geometry}
\usepackage{graphicx}
\usepackage{listings}
\usepackage{xcolor}
\usepackage{hyperref}
\usepackage{booktabs}
\usepackage{array}
\usepackage{longtable}
\usepackage{fancyhdr}
\usepackage{titlesec}
\usepackage{setspace}
\usepackage{parskip}
\usepackage{enumitem}
\usepackage{mdframed}
\usepackage{tikz}
\usetikzlibrary{shapes.geometric, arrows.meta, positioning, fit, backgrounds}

\hypersetup{
    colorlinks=true,
    linkcolor=blue!60!black,
    urlcolor=blue!60!black,
    citecolor=blue!60!black
}

% ── Code listings ─────────────────────────────────────────────────────────────
\definecolor{codebg}{RGB}{248,248,248}
\definecolor{codeframe}{RGB}{200,200,200}
\definecolor{keyword}{RGB}{0,0,180}
\definecolor{comment}{RGB}{80,120,80}
\definecolor{string}{RGB}{160,32,32}

\lstset{
    backgroundcolor=\color{codebg},
    frame=single,
    rulecolor=\color{codeframe},
    basicstyle=\ttfamily\footnotesize,
    keywordstyle=\color{keyword}\bfseries,
    commentstyle=\color{comment}\itshape,
    stringstyle=\color{string},
    numbers=left,
    numberstyle=\tiny\color{gray},
    numbersep=8pt,
    breaklines=true,
    captionpos=b,
    showstringspaces=false,
    tabsize=2
}

\lstdefinelanguage{YAML}{
  keywords={true,false,null,on,off},
  keywordstyle=\color{keyword}\bfseries,
  sensitive=false,
  comment=[l]{\#},
  commentstyle=\color{comment}\itshape,
  stringstyle=\color{string},
  morestring=[b]',
  morestring=[b]"
}

\lstdefinelanguage{Groovy}{
  keywords={pipeline, agent, environment, stages, stage, steps, post, when,
            branch, sh, parallel, always, success, failure, unstable,
            credentials, tools, options, def, if, else, return},
  keywordstyle=\color{keyword}\bfseries,
  sensitive=true,
  comment=[l]{//},
  commentstyle=\color{comment}\itshape,
  stringstyle=\color{string},
  morestring=[b]',
  morestring=[b]"
}

% ── En-têtes ──────────────────────────────────────────────────────────────────
\pagestyle{fancy}
\fancyhf{}
\fancyhead[L]{\small\leftmark}
\fancyhead[R]{\small BiblioGest CI/CD}
\fancyfoot[C]{\thepage}
\renewcommand{\headrulewidth}{0.4pt}

% ── Titres de chapitres ───────────────────────────────────────────────────────
\titleformat{\chapter}[block]
  {\normalfont\Large\bfseries\color{blue!50!black}}
  {\thechapter.}{0.8em}{}
\titlespacing*{\chapter}{0pt}{10pt}{20pt}

\onehalfspacing

% ══════════════════════════════════════════════════════════════════════════════
\begin{document}

% ── Page de titre ─────────────────────────────────────────────────────────────
\begin{titlepage}
\centering
\vspace*{1cm}
{\scshape\Large Rapport de Projet\par}
\vspace{0.3cm}
{\large Module~: Gestion des Projets IT\par}
\vspace{1.5cm}
{\Huge\bfseries Pipeline CI/CD Moderne Multi-Outils\par}
\vspace{0.5cm}
{\large Sujet 1~-- Application sur le projet BiblioGest\par}
\vspace{1.8cm}
\rule{\textwidth}{0.5pt}
\vspace{0.4cm}
{\large\bfseries GitHub Actions $\cdot$ GitLab CI/CD $\cdot$ Jenkins $\cdot$ CircleCI $\cdot$ Travis CI \\
Docker $\cdot$ SonarCloud / SonarQube $\cdot$ Trivy $\cdot$ Slack\par}
\vspace{0.4cm}
\rule{\textwidth}{0.5pt}
\vspace{2cm}

\begin{minipage}{0.45\textwidth}
\raggedright
\textbf{Filière~:} Génie Logiciel\\
\textbf{Module~:} Gestion des Projets IT
\end{minipage}
\hfill
\begin{minipage}{0.45\textwidth}
\raggedleft
\textbf{Année universitaire~:} 2025--2026
\end{minipage}

\vfill
{\large\today\par}
\end{titlepage}

% ── Résumé ────────────────────────────────────────────────────────────────────
\chapter*{Résumé}
\addcontentsline{toc}{chapter}{Résumé}

Ce rapport présente la mise en \oe{}uvre d'une pipeline complète d'intégration
et de déploiement continus (CI/CD) pour une application Jakarta~EE de gestion
de bibliothèques, BiblioGest. L'objectif est d'automatiser l'ensemble du cycle
de vie logiciel~: analyse statique (lint), tests unitaires et d'intégration,
analyse qualité, packaging, containerisation, scan de sécurité et notification.

La même logique est implémentée sur cinq moteurs CI/CD distincts -- \textbf{GitHub Actions},
\textbf{GitLab CI/CD}, \textbf{Jenkins}, \textbf{CircleCI} et \textbf{Travis CI} --
afin de comparer concrètement leurs syntaxes, leurs paradigmes et leurs
écosystèmes. La qualité du code est analysée par \textbf{SonarCloud} (mode SaaS)
ainsi que par une instance \textbf{SonarQube} auto-hébergée (mode on-premise),
les deux modes étant comparés en termes de coût, de maintenance et de
confidentialité. Le \textbf{Dockerfile multi-stage} produit une image runtime
allégée, l'image est scannée par \textbf{Trivy} (CVE), publiée simultanément
sur \textbf{Docker Hub} et \textbf{GitHub Container Registry}, et le résultat
final du pipeline est notifié à l'équipe via \textbf{Slack}.

Les tests, rédigés avec JUnit~5 et Mockito, couvrent la logique métier sans
aucune dépendance à la base de données Oracle, ce qui les rend exécutables en
quelques secondes sur n'importe quel runner CI/CD.

\tableofcontents
\listoffigures
\listoftables

% ══════════════════════════════════════════════════════════════════════════════
\chapter{Introduction}

\section{Contexte et motivation}

Le développement logiciel moderne fait face à une contradiction bien connue~:
plus un projet grandit, plus les tâches répétitives (compilation, exécution des
tests, déploiement) deviennent chronophages et sujettes aux erreurs humaines.
L'intégration continue (CI) et le déploiement continu (CD) constituent la
réponse industrielle à ce problème.

L'idée fondatrice remonte aux pratiques Extreme Programming des années~2000,
popularisées notamment par Martin Fowler. Elle consiste à fusionner
fréquemment les modifications de code dans une branche principale et à
déclencher automatiquement une série de vérifications à chaque intégration.
Lorsqu'on étend cette logique jusqu'à la mise en production automatique, on
parle de déploiement continu.

\section{Objectifs du projet}

Le sujet retenu pour le projet du module \emph{Gestion des Projets IT} est le
\textbf{Sujet 1~-- CI/CD moderne multi-outils}. Ses objectifs explicites sont~:

\begin{enumerate}[leftmargin=*, label=\arabic*.]
  \item Comparer les principaux moteurs CI/CD du marché~: déclencheurs,
        runners, jobs, stages, artefacts, secrets, notifications.
  \item Mettre en place une pipeline standard incluant lint, tests unitaires,
        tests d'intégration, build, analyse qualité et packaging Docker.
  \item Adapter la même logique sur plusieurs outils pour comprendre les
        différences de syntaxe et de fonctionnement.
  \item Développer une application support et automatiser tout son cycle.
  \item Aller plus loin avec~: publication d'image Docker sur un registre,
        intégration SonarQube, notification Slack/Teams.
\end{enumerate}

\section{Application support choisie}

L'application support est \textbf{BiblioGest}, un système de gestion de
bibliothèques développé en \textbf{Jakarta EE~10} (Servlets, JSP, JAX-RS),
packagé en WAR et déployé sur Apache Tomcat~10.1. La persistance s'appuie sur
Hibernate~ORM~6 et une base Oracle XE. L'application expose une API REST
\texttt{/api/livres} utilisée comme point de contrôle pour les healthchecks et
les tests d'intégration.

\section{Structure du rapport}

Le rapport s'organise comme suit~: le chapitre~\ref{ch:fondements} pose les
fondements théoriques du CI/CD et compare les cinq moteurs étudiés. Le
chapitre~\ref{ch:app} décrit l'application BiblioGest. Le chapitre~\ref{ch:archi}
détaille l'architecture de la pipeline implémentée. Les chapitres
\ref{ch:gha} à \ref{ch:travis} présentent l'implémentation sur chacun des cinq
moteurs. Le chapitre~\ref{ch:sonar} compare SonarCloud et SonarQube
self-hosted. Le chapitre~\ref{ch:diff} analyse les difficultés rencontrées.
Le chapitre~\ref{ch:reco} conclut par une comparaison finale et des
recommandations.

% ══════════════════════════════════════════════════════════════════════════════
\chapter{Fondements du CI/CD et panorama des outils}
\label{ch:fondements}

\section{Concepts communs}

Quelle que soit la plateforme utilisée, une pipeline CI/CD repose sur un
vocabulaire commun qu'il convient de définir avec précision avant toute
comparaison.

\begin{description}[leftmargin=2cm, style=nextline]
  \item[Déclencheur (trigger)] Événement qui lance l'exécution de la pipeline~:
        push sur une branche, ouverture d'une pull request, planification
        temporelle (cron), ou déclenchement manuel.
  \item[Runner / Agent] Machine (physique ou virtuelle) sur laquelle les jobs
        s'exécutent. Peut être hébergée par le fournisseur (hosted) ou
        autogérée (self-hosted).
  \item[Job] Unité d'exécution autonome composée d'une suite de steps. Les
        jobs peuvent être parallèles ou séquentiels selon leurs dépendances.
  \item[Stage] Groupe logique de jobs partageant la même phase du cycle de vie
        (lint, build, test, deploy). Concept central de GitLab et Travis,
        modélisé différemment dans les autres outils.
  \item[Artefact] Fichier produit par un job et transmis aux jobs suivants ou
        archivé~: rapports de tests, WAR, image Docker.
  \item[Secret] Variable sensible (mot de passe, token) injectée dans
        l'environnement d'exécution sans jamais apparaître dans les logs.
  \item[Cache] Répertoire conservé entre les exécutions pour accélérer les
        builds (dépendances Maven, layers Docker, modules Node).
  \item[Notification] Message émis à la fin du pipeline (succès ou échec) vers
        Slack, Teams, e-mail ou tout autre canal.
\end{description}

\section{Panorama des cinq outils étudiés}

\subsection{GitHub Actions}

Intégrée directement dans GitHub depuis~2019, GitHub~Actions adopte une
philosophie événementielle. Chaque workflow est un fichier~YAML stocké dans
\texttt{.github/workflows/}. La marketplace propose plus de~21\,000 actions
réutilisables couvrant la quasi-totalité des besoins courants.

Le modèle de facturation est généreux pour les projets publics~: les runners
Linux sont gratuits et illimités. Pour les projets privés, 2\,000~minutes
mensuelles sont offertes sur le plan gratuit. Les dépendances entre jobs sont
exprimées via \texttt{needs:}, ce qui autorise des graphes de dépendances
arbitraires (DAG), plus flexibles que la notion stricte de stages séquentiels.

\subsection{GitLab CI/CD}

GitLab~CI/CD est le système natif de la plateforme GitLab, défini dans un
fichier \texttt{.gitlab-ci.yml} à la racine du dépôt. Son concept central est la
notion de \emph{stages}~: tous les jobs d'un même stage s'exécutent en
parallèle, et le stage suivant ne démarre que si tous les jobs du précédent
réussissent (sauf \texttt{allow\_failure}).

GitLab propose des runners partagés (shared runners) dans sa version cloud
SaaS, mais l'hébergement self-managed reste très répandu en entreprise,
notamment pour des raisons de conformité des données. L'interface de
visualisation des pipelines est considérée comme l'une des meilleures du marché.

\subsection{Jenkins}

Jenkins est l'ancêtre des outils CI/CD modernes (premier release en~2011 sous
ce nom, héritier de Hudson). Open source, il s'installe on-premise et se
configure via une interface web ou, depuis l'introduction des Pipelines en~2016,
via un \texttt{Jenkinsfile} écrit en Groovy~DSL.

Sa force est son écosystème de plus de~1\,900 plugins, mais cela constitue
aussi sa faiblesse~: la gestion des dépendances entre plugins peut devenir un
véritable casse-tête de maintenance. Jenkins est le seul des cinq outils
étudiés à être \textbf{self-hosted par nature}~: il n'existe pas de mode SaaS
officiel.

\subsection{CircleCI}

CircleCI mise sur la performance et la simplicité de configuration. Son
concept distinctif est l'\textbf{orb}~: un module YAML réutilisable et
versionné, analogue aux actions GitHub mais avec une granularité plus large
(un orb peut contenir plusieurs jobs, commands et executors). La particularité
la plus notable est la possibilité de déboguer un build en SSH directement
dans le runner, ce qui s'avère précieux lors de problèmes d'environnement
difficiles à reproduire localement.

\subsection{Travis CI}

Travis~CI a longtemps été la référence pour les projets open source hébergés
sur GitHub. Son déclin relatif depuis~2020 (changement de politique tarifaire,
rachat par Idera, fuite vers GitHub Actions) a poussé une grande partie de la
communauté à migrer ailleurs. Il reste néanmoins utilisé dans de nombreux
projets existants et mérite d'être connu pour comprendre l'évolution
historique du domaine.

Travis se distingue par sa syntaxe particulièrement \textbf{concise}~: trois
lignes suffisent à déclencher un build Java standard, sans avoir à décrire
explicitement les étapes de checkout ou de configuration de JDK.

\section{Tableau comparatif initial}

\begin{table}[h!]
\centering
\small
\renewcommand{\arraystretch}{1.4}
\begin{tabular}{>{\bfseries}p{3cm} p{2.1cm} p{2.1cm} p{2.1cm} p{2.1cm} p{2.1cm}}
\toprule
Critère & GitHub Actions & GitLab CI/CD & Jenkins & CircleCI & Travis CI \\
\midrule
Config. fichier        & YAML & YAML & Groovy DSL & YAML & YAML \\
Hébergement            & SaaS & SaaS + SH & On-premise & SaaS & SaaS \\
Concept central        & workflows / jobs & stages / jobs & pipeline / stages & workflows / jobs & stages / jobs \\
Marketplace            & 21\,000+ actions & Modéré & 1\,900+ plugins & Orbs & Limité \\
Intégration Git        & Native GitHub & Native GitLab & Via plugin & GitHub/GitLab/BB & GitHub \\
Parallélisme           & Via \texttt{needs} & Par stage & Via agents & Natif (workflows) & Par stage \\
Docker natif           & Oui & Oui (DinD) & Plugin & Oui & Oui \\
Courbe d'apprentissage & Faible & Faible & Élevée & Moyenne & Faible \\
\bottomrule
\end{tabular}
\caption{Comparaison initiale des cinq outils CI/CD}
\label{tab:initial-comparison}
\end{table}

Une comparaison plus détaillée, enrichie par l'expérience d'implémentation,
est présentée au chapitre~\ref{ch:reco}.

% ══════════════════════════════════════════════════════════════════════════════
\chapter{Présentation de l'application BiblioGest}
\label{ch:app}

\section{Description fonctionnelle}

BiblioGest est une application web de gestion de bibliothèques. Elle permet à
des gestionnaires authentifiés d'administrer le catalogue de livres, les
auteurs, les catégories, les clients et les emprunts. Un seul administrateur
peut gérer les comptes utilisateurs.

Le modèle de données central est l'\textbf{emprunt}~: un client emprunte un
livre pour une durée déterminée, paie des frais calculés en fonction du tarif
journalier, et s'expose à une pénalité de~1,5\,× le tarif journalier par jour
de retard.

\section{Architecture technique}

L'application suit le pattern MVC Jakarta~EE classique~:

\begin{itemize}
  \item \textbf{Modèle}~: entités JPA/Hibernate~6 annotées
        (\texttt{@Entity}, \texttt{@ManyToMany}, etc.) et DAOs à transactions
        manuelles (\texttt{RESOURCE\_LOCAL}).
  \item \textbf{Contrôleur}~: Servlets annotées (\texttt{@WebServlet}) suivant
        le pattern Post-Redirect-Get.
  \item \textbf{Vue}~: pages JSP avec Tailwind CSS (CDN) et JSTL~3.0.
  \item \textbf{API REST}~: ressource JAX-RS \texttt{/api/livres} exposant le
        catalogue en JSON via Jersey~3.1.
\end{itemize}

\begin{figure}[h!]
\centering
\begin{tikzpicture}[
  box/.style={rectangle, draw=blue!40, fill=blue!5, rounded corners,
              minimum width=2.8cm, minimum height=1cm, text centered,
              font=\small},
  arr/.style={-Stealth, thick, blue!60}
]
\node[box] (client)  {Navigateur};
\node[box, right=1.8cm of client]  (servlet) {Servlet / JAX-RS};
\node[box, right=1.8cm of servlet] (dao)     {DAO / JPA};
\node[box, right=1.8cm of dao]     (db)      {Oracle XE};

\draw[arr] (client) -- (servlet) node[midway,above]{\tiny HTTP};
\draw[arr] (servlet) -- (dao) node[midway,above]{\tiny Appel};
\draw[arr] (dao) -- (db) node[midway,above]{\tiny JDBC};
\draw[arr, dashed] (db) -- (dao);
\draw[arr, dashed] (dao) -- (servlet);
\draw[arr, dashed] (servlet) -- (client) node[midway,below]{\tiny HTML/JSON};
\end{tikzpicture}
\caption{Architecture applicative de BiblioGest}
\label{fig:archi-app}
\end{figure}

\section{Route fonctionnelle exposée}

La route \texttt{GET /api/livres} est la route principale utilisée dans la
pipeline CI/CD~: elle est testée en intégration et sert de point de contrôle
pour le healthcheck Docker.

\begin{lstlisting}[language=Java, caption={Endpoint REST -- LivreResource.java}]
@GET
public Response getAll(@QueryParam("q") String q,
                       @QueryParam("disponible") Boolean disponible) {
    List<Livre> list = (q != null && !q.isBlank())
            ? dao.search(q) : dao.findAllWithRelations();
    if (Boolean.TRUE.equals(disponible))
        list = list.stream().filter(Livre::isDisponible).toList();
    return Response.ok(list.stream().map(this::toMap).toList()).build();
}
\end{lstlisting}

% ══════════════════════════════════════════════════════════════════════════════
\chapter{Architecture de la pipeline CI/CD}
\label{ch:archi}

\section{Vue d'ensemble}

La pipeline finale est composée de \textbf{sept étapes} exécutées dans l'ordre
suivant, identique sur les cinq moteurs CI/CD~:

\begin{figure}[h!]
\centering
\begin{tikzpicture}[
  stage/.style={rectangle, draw=blue!50, fill=blue!8, rounded corners=4pt,
                minimum width=2.1cm, minimum height=1.0cm, text centered,
                font=\scriptsize\bfseries},
  arr/.style={-Stealth, thick}
]
\node[stage] (l) {Lint};
\node[stage, right=0.6cm of l] (t) {Tests};
\node[stage, right=0.6cm of t] (q) {Qualité};
\node[stage, right=0.6cm of q] (p) {Package};
\node[stage, right=0.6cm of p] (d) {Docker};
\node[stage, right=0.6cm of d] (s) {Sécurité};
\node[stage, right=0.6cm of s] (n) {Notify};

\draw[arr] (l) -- (t);
\draw[arr] (t) -- (q);
\draw[arr] (t.south) to[bend right=20] (p.south);
\draw[arr] (p) -- (d);
\draw[arr] (d) -- (s);
\draw[arr] (s) -- (n);

\node[font=\tiny, below=0.3cm of q] {main};
\node[font=\tiny, below=0.3cm of d] {main};
\node[font=\tiny, below=0.3cm of s] {main};
\node[font=\tiny, below=0.3cm of n] {main};
\end{tikzpicture}
\caption{Architecture de la pipeline en sept étapes}
\label{fig:pipeline}
\end{figure}

\begin{table}[h!]
\centering
\small
\renewcommand{\arraystretch}{1.3}
\begin{tabular}{>{\bfseries}p{2.6cm} p{8.5cm} p{2cm}}
\toprule
Étape & Description & Branche \\
\midrule
Lint        & Checkstyle (style) + SpotBugs (sémantique) & toutes \\
Tests       & Unitaires (Surefire) + intégration (Failsafe) + couverture JaCoCo & toutes \\
Qualité     & Analyse SonarCloud avec quality gate & main \\
Package     & Build du WAR et upload de l'artefact & toutes \\
Docker      & Build multi-stage et push sur Docker Hub + GHCR & main \\
Sécurité    & Scan Trivy CVE de l'image Docker & main \\
Notify      & Notification Slack succès / échec & main \\
\bottomrule
\end{tabular}
\caption{Étapes de la pipeline et conditions d'exécution}
\label{tab:pipeline-stages}
\end{table}

\section{Stratégie de test sans base de données}

L'application utilise Oracle XE pour toutes ses opérations de persistance, ce
qui pose un problème d'environnement en CI~: on ne peut pas raisonnablement
lancer un conteneur Oracle XE lors de chaque exécution de la pipeline (temps
de démarrage de 90 secondes à 3 minutes, images de plusieurs Go). La solution
adoptée est de tester uniquement la logique métier indépendante de la couche
de persistance.

Deux classes ont été extraites à cet effet~:

\begin{itemize}
  \item \textbf{PasswordUtil}~: hachage et vérification SHA-256 des mots de
        passe. Pure Java, sans état.
  \item \textbf{BiblioUtils}~: calcul des frais d'emprunt et des pénalités de
        retard. Fonctions pures opérant sur \texttt{BigDecimal} et
        \texttt{LocalDate}.
\end{itemize}

Les tests d'intégration vérifient la couche REST sans démarrer de serveur~: le
DAO est remplacé par un mock Mockito injecté via le constructeur de
\texttt{LivreResource}.

\section{Convention de nommage}

\begin{itemize}
  \item Suffixe \texttt{Test} (ex.~\texttt{PasswordUtilTest})~: exécuté par
        Maven~Surefire pendant \texttt{mvn test}.
  \item Suffixe \texttt{IT} (ex.~\texttt{LivreResourceIT})~: exécuté par
        Maven~Failsafe pendant \texttt{mvn verify}.
\end{itemize}

Cette séparation explicite évite que les tests d'intégration ralentissent les
boucles de feedback rapides (tests unitaires en quelques secondes).

\section{Configuration du POM Maven}

Le \texttt{pom.xml} déclare cinq plugins clés~:

\begin{table}[h!]
\centering
\small
\renewcommand{\arraystretch}{1.3}
\begin{tabular}{>{\bfseries}p{4.5cm} p{2.5cm} p{6cm}}
\toprule
Plugin & Version & Rôle \\
\midrule
maven-surefire-plugin   & 3.2.5    & Exécution des tests unitaires \\
maven-failsafe-plugin   & 3.2.5    & Exécution des tests d'intégration \\
jacoco-maven-plugin     & 0.8.11   & Couverture de code (XML pour Sonar) \\
sonar-maven-plugin      & 3.11.0   & Envoi des résultats à SonarCloud \\
maven-checkstyle-plugin & 3.3.1    & Lint de style (Checkstyle 10.14.2) \\
spotbugs-maven-plugin   & 4.8.3.1  & Lint sémantique / détection de bugs \\
\bottomrule
\end{tabular}
\caption{Plugins Maven configurés dans le POM}
\label{tab:maven-plugins}
\end{table}

\section{Architecture du Dockerfile multi-stage}

Le Dockerfile est organisé en deux étapes (multi-stage build)~:

\begin{enumerate}
  \item \textbf{Stage \texttt{builder}}~: image \texttt{eclipse-temurin:17-jdk-alpine}
        avec Maven. Compile le projet et produit \texttt{target/bibliotheque.war}.
  \item \textbf{Stage \texttt{runtime}}~: image \texttt{tomcat:10.1-jdk17-temurin-alpine}.
        Copie uniquement le WAR depuis le builder, configure les variables
        d'environnement, expose le port~8080 et démarre Tomcat avec un
        utilisateur non-root.
\end{enumerate}

L'avantage est double~: l'image finale ne contient ni Maven ni le code source,
réduisant sa taille et sa surface d'attaque.

\begin{figure}[h!]
\centering
\begin{tikzpicture}[
  layer/.style={rectangle, draw=blue!50, fill=blue!8, rounded corners=4pt,
                minimum width=4cm, minimum height=0.9cm, text centered, font=\small},
  arr/.style={-Stealth, thick}
]
\node[layer] (src)     {Code source + pom.xml};
\node[layer, below=0.3cm of src, fill=blue!15] (builder)
  {\textbf{Stage builder} (JDK + Maven)};
\node[layer, below=0.3cm of builder] (war) {target/bibliotheque.war};
\node[layer, below=0.5cm of war, fill=green!15] (runtime)
  {\textbf{Stage runtime} (Tomcat seul)};
\node[layer, below=0.3cm of runtime] (img) {Image finale (\textasciitilde 250 Mo)};

\draw[arr] (src) -- (builder);
\draw[arr] (builder) -- (war);
\draw[arr] (war) -- (runtime) node[midway, right, font=\tiny] {COPY --from=builder};
\draw[arr] (runtime) -- (img);
\end{tikzpicture}
\caption{Architecture du Dockerfile multi-stage}
\label{fig:dockerfile-multistage}
\end{figure}

% ══════════════════════════════════════════════════════════════════════════════
\chapter{Implémentation sur GitHub Actions}
\label{ch:gha}

\section{Fichier de configuration}

Le workflow est défini dans \texttt{.github/workflows/ci-cd.yml}. Sept jobs y
sont déclarés~: \texttt{lint}, \texttt{test}, \texttt{quality}, \texttt{package},
\texttt{docker}, \texttt{security-scan}, \texttt{notify}.

\section{Déclencheurs}

\begin{lstlisting}[language=YAML, caption={Déclencheurs du workflow}]
on:
  push:
    branches: [main, develop]
  pull_request:
    branches: [main]
\end{lstlisting}

Tout push sur \texttt{main} ou \texttt{develop} déclenche les jobs
\texttt{lint}, \texttt{test} et \texttt{package}. Les jobs \texttt{quality},
\texttt{docker}, \texttt{security-scan} et \texttt{notify} sont conditionnés
par \texttt{if:\ github.ref == 'refs/heads/main'} pour éviter de consommer le
quota SonarCloud et de pousser des images intermédiaires sur Docker Hub à
chaque pull request.

\section{Job lint}

\begin{lstlisting}[language=YAML, caption={Job lint -- Checkstyle + SpotBugs}]
lint:
  name: Lint (Checkstyle + SpotBugs)
  runs-on: ubuntu-latest
  steps:
    - uses: actions/checkout@v4
    - uses: actions/setup-java@v4
      with:
        java-version: '17'
        distribution: 'temurin'
        cache: maven
    - run: mvn -B checkstyle:check --no-transfer-progress
    - run: mvn -B compile -DskipTests --no-transfer-progress
    - run: mvn -B spotbugs:check --no-transfer-progress
      continue-on-error: true
    - if: always()
      uses: actions/upload-artifact@v4
      with:
        name: lint-reports
        path: |
          target/checkstyle-result.xml
          target/spotbugsXml.xml
\end{lstlisting}

Note pédagogique~: Checkstyle et SpotBugs sont configurés en mode
\textbf{informational}, c'est-à-dire que les violations sont rapportées
(rapports XML uploadés en artefacts) sans bloquer le pipeline. Ce choix
s'explique par le fait que le code applicatif existant n'a pas été écrit avec
ces outils en tête (177 violations Checkstyle initialement). Rendre le job
bloquant aurait nécessité une réécriture massive sortant du cadre du sujet.

\section{Job test}

\begin{lstlisting}[language=YAML, caption={Job test -- extrait}]
test:
  name: Tests
  runs-on: ubuntu-latest
  needs: lint
  steps:
    - uses: actions/checkout@v4
      with:
        fetch-depth: 0
    - uses: actions/setup-java@v4
      with:
        java-version: '17'
        distribution: 'temurin'
        cache: maven
    - run: mvn test -B --no-transfer-progress
    - run: mvn verify -B --no-transfer-progress
    - if: always()
      uses: actions/upload-artifact@v4
      with:
        name: test-reports
        path: |
          target/surefire-reports/
          target/failsafe-reports/
          target/site/jacoco/
\end{lstlisting}

Le cache Maven est activé via l'option \texttt{cache: maven} de
\texttt{actions/setup-java}. Il réduit significativement la durée des builds
en conservant les dépendances téléchargées entre les exécutions (de
\textasciitilde 60s à \textasciitilde 8s pour la phase de résolution Maven).

\section{Job docker -- publication multi-registry}

L'image est publiée simultanément sur \textbf{Docker Hub} (public) et sur
\textbf{GitHub Container Registry} (GHCR, intégré à GitHub) grâce à l'action
\texttt{docker/metadata-action} qui génère plusieurs tags à la volée~:

\begin{lstlisting}[language=YAML, caption={Job docker -- multi-registry}]
docker:
  name: Docker Build & Push
  runs-on: ubuntu-latest
  needs: package
  if: github.ref == 'refs/heads/main'
  permissions:
    contents: read
    packages: write
  steps:
    - uses: actions/checkout@v4
    - uses: docker/setup-buildx-action@v3
    - uses: docker/login-action@v3
      with:
        username: ${{ secrets.DOCKER_USERNAME }}
        password: ${{ secrets.DOCKER_PASSWORD }}
    - uses: docker/login-action@v3
      with:
        registry: ghcr.io
        username: ${{ github.actor }}
        password: ${{ secrets.GITHUB_TOKEN }}
    - id: meta
      uses: docker/metadata-action@v5
      with:
        images: |
          ${{ secrets.DOCKER_USERNAME }}/bibliotheque
          ghcr.io/${{ github.repository_owner }}/bibliotheque
        tags: |
          type=raw,value=latest,enable=true
          type=sha,prefix=,format=long
    - uses: docker/build-push-action@v5
      with:
        context: .
        push: true
        tags: ${{ steps.meta.outputs.tags }}
        labels: ${{ steps.meta.outputs.labels }}
        cache-from: type=gha
        cache-to: type=gha,mode=max
\end{lstlisting}

\section{Job security-scan -- analyse Trivy}

Trivy scanne l'image publiée à la recherche de CVE (Common Vulnerabilities and
Exposures). Les résultats sont uploadés au format SARIF dans l'onglet
\textbf{Security} de GitHub.

\begin{lstlisting}[language=YAML, caption={Job security-scan -- Trivy}]
security-scan:
  name: Scan securite Trivy
  runs-on: ubuntu-latest
  needs: docker
  if: github.ref == 'refs/heads/main'
  steps:
    - uses: aquasecurity/trivy-action@0.24.0
      with:
        image-ref: ${{ secrets.DOCKER_USERNAME }}/bibliotheque:latest
        format: 'table'
        severity: 'CRITICAL,HIGH'
        exit-code: '0'
        ignore-unfixed: true
    - uses: aquasecurity/trivy-action@0.24.0
      with:
        image-ref: ${{ secrets.DOCKER_USERNAME }}/bibliotheque:latest
        format: 'sarif'
        output: 'trivy-results.sarif'
    - uses: github/codeql-action/upload-sarif@v3
      with:
        sarif_file: 'trivy-results.sarif'
\end{lstlisting}

\section{Job notify -- notification Slack}

Le dernier job s'exécute \emph{toujours} (\texttt{if: always()}) pour notifier
aussi bien les succès que les échecs~:

\begin{lstlisting}[language=YAML, caption={Job notify -- Slack}]
notify:
  needs: [lint, test, quality, package, docker, security-scan]
  if: always() && github.ref == 'refs/heads/main'
  runs-on: ubuntu-latest
  steps:
    - id: status
      run: |
        if [ "${{ contains(needs.*.result, 'failure') }}" = "true" ]; then
          echo "text=Pipeline ECHOUEE" >> $GITHUB_OUTPUT
          echo "color=danger" >> $GITHUB_OUTPUT
        else
          echo "text=Pipeline REUSSIE" >> $GITHUB_OUTPUT
          echo "color=good" >> $GITHUB_OUTPUT
        fi
    - uses: slackapi/slack-github-action@v1.27.0
      with:
        payload: |
          {
            "attachments": [{
              "color": "${{ steps.status.outputs.color }}",
              "blocks": [...]
            }]
          }
      env:
        SLACK_WEBHOOK_URL: ${{ secrets.SLACK_WEBHOOK_URL }}
        SLACK_WEBHOOK_TYPE: INCOMING_WEBHOOK
\end{lstlisting}

\section{Secrets et variables}

\begin{table}[h!]
\centering
\small
\renewcommand{\arraystretch}{1.3}
\begin{tabular}{>{\bfseries}p{4.5cm} p{1.8cm} p{6cm}}
\toprule
Nom & Type & Description \\
\midrule
DOCKER\_USERNAME    & Secret   & Identifiant Docker Hub \\
DOCKER\_PASSWORD    & Secret   & Token Docker Hub (PAT) \\
SONAR\_TOKEN        & Secret   & Token SonarCloud \\
SLACK\_WEBHOOK\_URL & Secret   & URL d'un Incoming Webhook Slack \\
SONAR\_PROJECT\_KEY & Variable & Clé du projet SonarCloud \\
SONAR\_ORGANIZATION & Variable & Organisation SonarCloud \\
\bottomrule
\end{tabular}
\caption{Secrets et variables GitHub Actions}
\end{table}

% ══════════════════════════════════════════════════════════════════════════════
\chapter{Implémentation sur GitLab CI/CD}
\label{ch:gitlab}

\section{Structure en stages}

Le fichier \texttt{.gitlab-ci.yml} déclare explicitement les sept stages~:

\begin{lstlisting}[language=YAML, caption={Stages GitLab}]
stages:
  - lint
  - test
  - quality
  - package
  - docker
  - security
  - notify
\end{lstlisting}

Tous les jobs d'un même stage s'exécutent en parallèle. Ainsi, le stage
\texttt{lint} contient deux jobs (\texttt{lint:checkstyle} et
\texttt{lint:spotbugs}) qui tournent simultanément sur des runners différents.

\section{Cache Maven mutualisé}

GitLab gère le cache via une clé partageable entre jobs~:

\begin{lstlisting}[language=YAML, caption={Cache Maven GitLab}]
.maven-cache: &maven-cache
  cache:
    key: "$CI_COMMIT_REF_SLUG-maven"
    paths:
      - .m2/repository
\end{lstlisting}

L'\textbf{anchor YAML} (\texttt{\&maven-cache}) est ensuite réutilisé dans chaque
job par \texttt{<<: *maven-cache}, évitant la duplication.

\section{Docker-in-Docker (DinD)}

Pour construire et pousser une image Docker, GitLab utilise le service DinD~:

\begin{lstlisting}[language=YAML, caption={Job docker:build-push}]
docker:build-push:
  stage: docker
  image: docker:24
  services:
    - docker:24-dind
  rules:
    - if: '$CI_COMMIT_BRANCH == "main"'
  needs:
    - job: package:war
      artifacts: true
  variables:
    DOCKER_TLS_CERTDIR: "/certs"
  before_script:
    - echo "$DOCKER_PASSWORD" | docker login -u "$DOCKER_USERNAME" --password-stdin
  script:
    - docker build -t $DOCKER_HUB_IMAGE:$CI_COMMIT_SHA -t $DOCKER_HUB_IMAGE:latest .
    - docker push $DOCKER_HUB_IMAGE:$CI_COMMIT_SHA
    - docker push $DOCKER_HUB_IMAGE:latest
\end{lstlisting}

\section{Intégration native JUnit}

GitLab affiche directement les résultats JUnit dans l'interface des Merge
Requests grâce à la directive \texttt{artifacts:reports:junit}~:

\begin{lstlisting}[language=YAML, caption={Reports JUnit GitLab}]
test:
  stage: test
  script:
    - mvn $MAVEN_CLI_OPTS verify
  artifacts:
    reports:
      junit:
        - target/surefire-reports/TEST-*.xml
        - target/failsafe-reports/TEST-*.xml
  coverage: '/Total.*?([0-9]{1,3})%/'
\end{lstlisting}

La regex \texttt{coverage} permet d'extraire le pourcentage de couverture des
logs et de l'afficher comme badge automatique sur le projet.

% ══════════════════════════════════════════════════════════════════════════════
\chapter{Implémentation sur Jenkins}
\label{ch:jenkins}

\section{Pipeline déclarative}

Le \texttt{Jenkinsfile} utilise la syntaxe \textbf{déclarative} (introduite en
2017), plus structurée et lisible que la syntaxe scriptée d'origine~:

\begin{lstlisting}[language=Groovy, caption={Squelette du Jenkinsfile}]
pipeline {
    agent any
    options {
        timeout(time: 30, unit: 'MINUTES')
        timestamps()
        disableConcurrentBuilds()
        buildDiscarder(logRotator(numToKeepStr: '10'))
    }
    environment {
        SONAR_TOKEN   = credentials('sonar-token')
        DOCKER_CREDS  = credentials('docker-hub')
        SLACK_WEBHOOK = credentials('slack-webhook')
    }
    tools {
        jdk   'jdk-17'
        maven 'maven-3.9'
    }
    stages { ... }
    post   { ... }
}
\end{lstlisting}

\section{Parallélisme de stages}

Jenkins permet de paralléliser plusieurs sous-stages dans un même bloc~:

\begin{lstlisting}[language=Groovy, caption={Lint parallèle dans Jenkins}]
stage('Lint') {
    parallel {
        stage('Checkstyle') {
            steps { sh 'mvn -B checkstyle:check' }
            post { always {
                recordIssues(tools: [checkStyle(pattern: 'target/checkstyle-result.xml')])
            }}
        }
        stage('SpotBugs') {
            steps {
                sh 'mvn -B compile -DskipTests'
                sh 'mvn -B spotbugs:check || true'
            }
            post { always {
                recordIssues(tools: [spotBugs(pattern: 'target/spotbugsXml.xml')])
            }}
        }
    }
}
\end{lstlisting}

\section{Gestionnaire de credentials}

Jenkins gère les secrets via son propre vault. Trois credentials doivent être
créés dans \textbf{Manage Jenkins~$\to$~Credentials}~:

\begin{table}[h!]
\centering
\small
\begin{tabular}{>{\bfseries}p{3cm} p{4cm} p{5cm}}
\toprule
ID & Type & Description \\
\midrule
sonar-token   & Secret text          & Token SonarCloud \\
docker-hub    & Username with password & Compte Docker Hub \\
slack-webhook & Secret text          & URL Slack webhook \\
\bottomrule
\end{tabular}
\caption{Credentials Jenkins requis}
\end{table}

L'usage en pipeline est~: \texttt{credentials('sonar-token')} expose la valeur
dans la variable d'environnement \texttt{SONAR\_TOKEN}, masquée dans tous les
logs.

\section{Post-actions et notification}

Le bloc \texttt{post} permet de déclencher des actions selon le résultat~:

\begin{lstlisting}[language=Groovy, caption={Post-actions Jenkins}]
post {
    success { slackNotify('good',   'Pipeline REUSSIE') }
    failure { slackNotify('danger', 'Pipeline ECHOUEE') }
    unstable { slackNotify('warning', 'Pipeline INSTABLE') }
    always  { cleanWs() }
}
\end{lstlisting}

\section{Plugins Jenkins requis}

\begin{itemize}
  \item \textbf{Pipeline}~: support des Jenkinsfile.
  \item \textbf{Docker Pipeline}~: exécution dans des conteneurs Docker.
  \item \textbf{Warnings Next Generation}~: agrégation des rapports Checkstyle / SpotBugs.
  \item \textbf{HTML Publisher}~: publication des rapports JaCoCo.
  \item \textbf{JUnit}~: collecte des résultats de tests.
  \item \textbf{Email Extension}~: envoi de mails enrichis.
  \item \textbf{Workspace Cleanup}~: nettoyage du workspace.
\end{itemize}

% ══════════════════════════════════════════════════════════════════════════════
\chapter{Implémentation sur CircleCI}
\label{ch:circleci}

\section{Orbs et executors}

CircleCI introduit deux concepts originaux par rapport aux autres outils~:

\begin{description}[leftmargin=1cm]
  \item[Orb] Module YAML versionné regroupant jobs, commands et executors
        réutilisables. Analogue à une action GitHub mais avec une granularité
        plus large.
  \item[Executor] Environnement d'exécution préconfiguré (Docker, machine VM,
        macOS) référençable par nom dans plusieurs jobs.
\end{description}

\begin{lstlisting}[language=YAML, caption={Orbs et executors CircleCI}]
version: 2.1
orbs:
  maven:  circleci/maven@1.4.1
  docker: circleci/docker@2.6.0

executors:
  java-executor:
    docker:
      - image: cimg/openjdk:17.0
    working_directory: ~/bibliotheque
    resource_class: medium

  docker-executor:
    machine:
      image: ubuntu-2204:current
    working_directory: ~/bibliotheque
\end{lstlisting}

\section{Workflow et filters}

Les dépendances entre jobs et les conditions de branche sont définies dans le
bloc \texttt{workflows}~:

\begin{lstlisting}[language=YAML, caption={Workflow CircleCI}]
workflows:
  bibliogest-pipeline:
    jobs:
      - lint
      - test:
          requires: [lint]
      - quality:
          requires: [test]
          filters:
            branches:
              only: main
      - package:
          requires: [test]
      - docker-build-push:
          requires: [package]
          filters:
            branches:
              only: main
\end{lstlisting}

\section{Workspaces}

Le partage d'artefacts entre jobs utilise \texttt{persist\_to\_workspace} et
\texttt{attach\_workspace}~:

\begin{lstlisting}[language=YAML, caption={Workspace CircleCI}]
- persist_to_workspace:
    root: .
    paths:
      - target/bibliotheque.war
      - Dockerfile
\end{lstlisting}

% ══════════════════════════════════════════════════════════════════════════════
\chapter{Implémentation sur Travis CI}
\label{ch:travis}

\section{Syntaxe minimaliste}

Travis CI est le plus concis des cinq outils. Une configuration Java basique
tient en trois lignes~:

\begin{lstlisting}[language=YAML, caption={Configuration Travis minimale}]
language: java
jdk: openjdk17
script: mvn verify
\end{lstlisting}

Le checkout, la configuration du JDK et le \texttt{mvn install} de base sont
implicites.

\section{Configuration avancée par stages}

Pour un pipeline plus complet, Travis utilise \texttt{stages} et \texttt{jobs.include}~:

\begin{lstlisting}[language=YAML, caption={Stages Travis}]
stages:
  - lint
  - test
  - name: quality
    if: branch = main
  - package
  - name: docker
    if: branch = main

jobs:
  include:
    - stage: lint
      script:
        - mvn -B checkstyle:check
        - mvn -B compile -DskipTests
        - mvn -B spotbugs:check || true
    - stage: test
      script:
        - mvn -B verify --no-transfer-progress
\end{lstlisting}

\section{Notifications natives}

Travis dispose d'un module de notification natif par e-mail~:

\begin{lstlisting}[language=YAML, caption={Notifications email Travis}]
notifications:
  email:
    on_success: change
    on_failure: always
\end{lstlisting}

L'option \texttt{change} signifie qu'un mail n'est envoyé qu'en cas de
changement d'état (succès après échec, ou échec après succès), évitant le spam
sur les builds successifs réussis.

\section{Avertissement}

Depuis 2021, Travis CI a changé sa politique tarifaire et n'est plus gratuit
même pour les projets OSS au-delà d'un quota initial. La communauté a
massivement migré vers GitHub Actions. Travis reste pédagogiquement intéressant
mais n'est plus recommandé pour de nouveaux projets.

% ══════════════════════════════════════════════════════════════════════════════
\chapter{SonarCloud vs SonarQube self-hosted}
\label{ch:sonar}

\section{Analyse de qualité}

L'analyse qualité repose sur le moteur \textbf{Sonar}, qui détecte~:

\begin{itemize}
  \item \textbf{Bugs}~: code potentiellement défaillant à l'exécution.
  \item \textbf{Vulnérabilités}~: failles de sécurité (CWE, OWASP).
  \item \textbf{Code smells}~: défauts de maintenabilité.
  \item \textbf{Hot spots de sécurité}~: code à risque méritant une revue.
  \item \textbf{Couverture}~: pourcentage de lignes exécutées par les tests
        (via JaCoCo).
  \item \textbf{Duplication}~: blocs de code similaires.
\end{itemize}

Deux modes de déploiement sont possibles~: SonarCloud (SaaS) et SonarQube
(auto-hébergé).

\section{Mode SaaS -- SonarCloud}

SonarCloud est la version hébergée dans le cloud de SonarQube. Configuration
dans toutes les pipelines~:

\begin{lstlisting}[language=bash, caption={Lancement de SonarCloud}]
mvn verify sonar:sonar \
  -Dsonar.projectKey=${SONAR_PROJECT_KEY} \
  -Dsonar.organization=${SONAR_ORGANIZATION} \
  -Dsonar.host.url=https://sonarcloud.io \
  -Dsonar.token=${SONAR_TOKEN}
\end{lstlisting}

\section{Mode on-premise -- SonarQube self-hosted}

Pour démonstration, un fichier \texttt{docker-compose.sonarqube.yml} déploie
une instance SonarQube complète avec sa base PostgreSQL~:

\begin{lstlisting}[language=YAML, caption={docker-compose.sonarqube.yml -- extrait}]
services:
  sonarqube-db:
    image: postgres:15-alpine
    environment:
      POSTGRES_USER: sonar
      POSTGRES_PASSWORD: sonar
      POSTGRES_DB: sonarqube
    volumes:
      - sonarqube-db-data:/var/lib/postgresql/data

  sonarqube:
    image: sonarqube:10.4-community
    depends_on:
      sonarqube-db:
        condition: service_healthy
    environment:
      SONAR_JDBC_URL: jdbc:postgresql://sonarqube-db:5432/sonarqube
    ports:
      - "9000:9000"
\end{lstlisting}

Démarrage~: \texttt{docker compose -f docker-compose.sonarqube.yml up -d}

Accès~: \url{http://localhost:9000} (admin / admin au premier login).

Analyse pointant sur l'instance locale~:

\begin{lstlisting}[language=bash]
mvn verify sonar:sonar \
  -Dsonar.host.url=http://localhost:9000 \
  -Dsonar.token=<token_genere>
\end{lstlisting}

\section{Comparaison synthétique}

\begin{table}[h!]
\centering
\small
\renewcommand{\arraystretch}{1.4}
\begin{tabular}{>{\bfseries}p{3.5cm} p{4.5cm} p{4.5cm}}
\toprule
Critère & SonarCloud & SonarQube self-hosted \\
\midrule
Hébergement              & SaaS (cloud Sonar) & À votre charge \\
Coût                     & Gratuit pour OSS public ; payant privé & Édition Community gratuite (autres éditions payantes) \\
Maintenance              & Aucune & Mises à jour, sauvegardes, sécurité \\
Confidentialité          & Code envoyé chez Sonar & 100\% on-prem \\
Setup CI/CD              & Token + URL & Token + URL exposée publiquement \\
Quality gates            & Préconfigurées + customisables & Customisables \\
Idéal pour               & Démos, OSS, startups & Entreprises, secteurs régulés \\
Effort initial           & Quelques minutes & Plusieurs heures (infra) \\
\bottomrule
\end{tabular}
\caption{Comparaison SonarCloud vs SonarQube self-hosted}
\label{tab:sonar-comparison}
\end{table}

% ══════════════════════════════════════════════════════════════════════════════
\chapter{Difficultés rencontrées}
\label{ch:diff}

\section{Tests sans base de données}

La première difficulté a concerné l'architecture des tests. L'application
utilise JPA avec Hibernate et une base Oracle XE. Le bloc statique de
\texttt{JpaUtil} tente de créer l'\texttt{EntityManagerFactory} dès le
chargement de la classe, ce qui échoue immédiatement si aucune base de données
n'est joignable.

La solution a consisté à n'instancier aucun DAO dans les tests~: les tests
unitaires testent des classes utilitaires pures (\texttt{PasswordUtil},
\texttt{BiblioUtils}), et les tests d'intégration injectent des mocks Mockito
dans \texttt{LivreResource} via un constructeur dédié.

\section{Constructeur d'injection package-private}

Un bug subtil pré-existait dans le projet~: le constructeur d'injection de
\texttt{LivreResource} était déclaré sans modificateur (donc package-private),
alors que la classe de test \texttt{LivreResourceIT} se trouve dans le package
\texttt{ma.bibliotheque.integration}. Cela provoquait une erreur de
compilation lors de \texttt{mvn test}.

Le bug n'avait jamais été détecté car le test n'avait apparemment jamais été
exécuté localement avant la mise en place du pipeline. La correction a
consisté à rendre le constructeur \texttt{public}, ce qui est légitime
puisqu'il sert à l'injection de dépendances.

\begin{lstlisting}[language=Java, caption={Correction du constructeur}]
// Avant (package-private, inaccessible depuis le test)
LivreResource(LivreDao dao) { this.dao = dao; }

// Après
public LivreResource(LivreDao dao) { this.dao = dao; }
\end{lstlisting}

\section{Conflit de jar dans WEB-INF/lib}

Lors du premier déploiement sur Tomcat, une \texttt{NoClassDefFoundError} sur
\texttt{jakarta/servlet/jsp/jstl/core/ConditionalTagSupport} s'est manifestée.
L'investigation a révélé que la dépendance Maven
\texttt{org.glassfish.web:jakarta.servlet.jsp.jstl:3.0.1} déclare l'API JSTL
en scope \texttt{provided} dans son propre POM, empêchant Maven de la résoudre
transitivement. La correction a consisté à déclarer explicitement
\texttt{jakarta.servlet.jsp.jstl-api:3.0.0} avec le scope \texttt{compile}.

Un second conflit concernait \texttt{org.glassfish:jakarta.el:4.0.2}~: ce jar
présent dans \texttt{WEB-INF/lib} entrait en conflit avec l'implémentation
EL~5.0 fournie par Tomcat~10.1, provoquant des \texttt{ClassCastException} au
bootstrap JSTL. La solution a été de passer ce jar en scope \texttt{provided}.

\section{Checkstyle bloquant sur code legacy}

À la première exécution, Checkstyle a remonté 177 violations sur le code
applicatif existant. Plusieurs choix étaient possibles~:

\begin{enumerate}
  \item Refactorer le code source pour passer Checkstyle (effort considérable,
        hors périmètre du sujet).
  \item Adopter une configuration Checkstyle ultra-permissive (vidant le sens
        de l'outil).
  \item Configurer le plugin Maven en mode \emph{informational}, c'est-à-dire
        que les violations sont rapportées sans bloquer le build.
\end{enumerate}

L'option 3 a été retenue~: le plugin est configuré avec \texttt{<failOnError>false</failOnError>}
et \texttt{<failOnViolation>false</failOnViolation>}, et les rapports XML sont
uploadés comme artefacts CI. Cette approche est cohérente avec le but
pédagogique du sujet (démontrer l'intégration de lint dans le pipeline) sans
pénaliser le projet existant.

\section{Oracle XE et Docker Compose}

L'image Docker d'Oracle XE prend entre 90 et 180~secondes à s'initialiser
complètement après le démarrage du conteneur. Sans la directive
\texttt{condition: service\_healthy} dans \texttt{docker-compose.yml},
l'application tente de se connecter avant que la PDB \texttt{XEPDB1} ne soit
disponible, ce qui déclenche une \texttt{ExceptionInInitializerError} non
récupérable dans \texttt{JpaUtil}. La configuration d'un healthcheck adéquat
sur le conteneur Oracle XE et la directive \texttt{depends\_on.condition} ont
été indispensables.

\section{Permissions GitHub Container Registry}

Pour pousser une image sur GHCR depuis GitHub Actions, le job doit déclarer
explicitement \texttt{permissions: packages: write}~:

\begin{lstlisting}[language=YAML]
docker:
  permissions:
    contents: read
    packages: write
\end{lstlisting}

Sans cela, l'authentification fonctionne mais le push échoue avec une erreur
403 difficile à diagnostiquer car les logs n'indiquent pas explicitement le
problème de permissions.

\section{SpotBugs et entités JPA}

SpotBugs détecte que les entités JPA exposent des collections internes via les
getters / setters (règle \texttt{EI\_EXPOSE\_REP}). Bien que ces avertissements
soient légitimes en Java pur, ils sont structurellement nécessaires pour
permettre à Hibernate de gérer les relations. La solution a été de configurer
SpotBugs avec \texttt{continue-on-error: true} côté CI et de documenter ces
warnings comme acceptables pour ce contexte JPA.

% ══════════════════════════════════════════════════════════════════════════════
\chapter{Comparaison finale et recommandations}
\label{ch:reco}

\section{Synthèse comparative détaillée}

Après implémentation effective sur les cinq moteurs, on peut produire une
comparaison fondée sur l'expérience plutôt que sur la documentation.

\begin{table}[h!]
\centering
\scriptsize
\renewcommand{\arraystretch}{1.4}
\begin{tabular}{>{\bfseries}p{3.2cm} p{1.9cm} p{1.9cm} p{1.9cm} p{1.9cm} p{1.9cm}}
\toprule
Critère & GitHub Actions & GitLab CI/CD & Jenkins & CircleCI & Travis CI \\
\midrule
Type             & SaaS & SaaS + SH & Self-host & SaaS & SaaS \\
Format config    & YAML & YAML & Groovy DSL & YAML & YAML \\
Concept central  & jobs / steps & stages / jobs & pipeline / stages & workflows / jobs & stages / jobs \\
Runner           & GitHub-hosted ou SH & GitLab Runner & Agents & CircleCI executors & Travis VM \\
Secrets          & Repo / org secrets & CI/CD variables & Credentials plugin & Project env vars & Encrypted env vars \\
Cache deps       & \texttt{actions/cache} & \texttt{cache:} & Plugin & \texttt{save\_cache} & \texttt{cache:} \\
Marketplace      & 21k+ actions & Modéré & 1.9k+ plugins & Orbs & Limité \\
Docker natif     & Oui (Buildx) & Oui (DinD) & Plugin & Oui (machine) & Oui \\
Parallélisme     & Via \texttt{needs} & Par stage & \texttt{parallel} & Natif & Par stage \\
Notification     & Action Slack & curl ou intégration & \texttt{post} block & Step Slack & \texttt{notifications:} \\
Gratuit OSS      & Oui (généreux) & Oui (limité) & Gratuit (self-host) & Limité & Payant 2021+ \\
Courbe apprent.  & Faible & Faible & Élevée & Moyenne & Faible \\
Vendor lock-in   & Faible & Faible & Nul & Moyen & Moyen \\
\bottomrule
\end{tabular}
\caption{Comparaison détaillée des cinq moteurs CI/CD}
\label{tab:final-comparison}
\end{table}

\section{Équivalences conceptuelles}

\begin{table}[h!]
\centering
\scriptsize
\renewcommand{\arraystretch}{1.3}
\begin{tabular}{>{\bfseries}p{2.5cm} p{2.5cm} p{2cm} p{2.5cm} p{2.3cm} p{2cm}}
\toprule
Concept & GitHub Actions & GitLab CI & Jenkins & CircleCI & Travis \\
\midrule
Unité d'exécution & job             & job        & stage         & job        & job \\
Dépendance        & \texttt{needs:} & \texttt{needs:} / stage & ordre dans \texttt{stages} & \texttt{requires:} & ordre dans \texttt{stages} \\
Condition         & \texttt{if:}    & \texttt{rules:} / \texttt{only:} & \texttt{when \{\}} & \texttt{filters:} & \texttt{if:} \\
Artefact          & upload-artifact & \texttt{artifacts:} & \texttt{archiveArtifacts} & \texttt{persist\_to\_workspace} & \texttt{workspaces:} \\
Cache             & actions/cache   & \texttt{cache:} & plugin & \texttt{save\_cache} & \texttt{cache:} \\
Secret            & \verb|${{ secrets.X }}| & \verb|$X| (masked) & \verb|credentials('id')| & \verb|$X| (project) & \verb|$X| (encrypted) \\
\bottomrule
\end{tabular}
\caption{Équivalences conceptuelles entre les cinq moteurs}
\label{tab:equivalences}
\end{table}

\section{Forces et limites de chaque outil}

\subsection{GitHub Actions}

\textbf{Forces}~: marketplace gigantesque, intégration GitHub sans friction,
runners gratuits illimités pour OSS, syntaxe YAML très lisible, support
matrix-build natif, intégration native avec GitHub Container Registry.

\textbf{Limites}~: dépendant de l'écosystème GitHub (vendor partiel), self-hosted
runners moins matures que GitLab Runner, debug en SSH non disponible.

\subsection{GitLab CI/CD}

\textbf{Forces}~: même outil pour SCM, CI/CD, registry, monitoring (DevOps
complet), interface de visualisation excellente, GitLab Runner très flexible,
mode self-hosted mature.

\textbf{Limites}~: marketplace plus modeste qu'Actions, syntaxe \texttt{rules:}
complexe pour les conditions avancées, performances des shared runners en
plan gratuit limitées.

\subsection{Jenkins}

\textbf{Forces}~: 100\% open source et auto-hébergé, écosystème de plugins
inégalé, flexibilité maximale (n'importe quel langage scriptable),
parfaitement intégré aux outils legacy d'entreprise.

\textbf{Limites}~: courbe d'apprentissage élevée (Groovy DSL), maintenance
lourde (mises à jour de plugins fréquentes, conflits possibles), pas de mode
SaaS officiel, UI vieillissante.

\subsection{CircleCI}

\textbf{Forces}~: performance élevée, orbs très bien conçus, debugging via
SSH dans le runner, support multi-VCS (GitHub, GitLab, Bitbucket).

\textbf{Limites}~: plan gratuit limité (6\,000 minutes/mois), moins d'intégrations
natives que GitHub Actions, syntaxe légèrement plus verbeuse.

\subsection{Travis CI}

\textbf{Forces}~: syntaxe minimaliste idéale pour démarrer, configuration de
base ultra-rapide, historiquement le standard OSS.

\textbf{Limites}~: déclin commercial depuis 2021, communauté en fuite vers
GitHub Actions, plan gratuit OSS supprimé, écosystème stagnant.

\section{Recommandations par contexte}

\begin{description}[leftmargin=0cm]
  \item[Projet académique ou OSS public hébergé sur GitHub]
        GitHub Actions sans hésitation~: gratuit, immédiat, marketplace
        couvrant 100\% des besoins.
  \item[Équipe utilisant GitLab self-hosted]
        GitLab CI/CD~: cohérence d'outillage, DevOps unifié, pas de
        synchronisation à maintenir.
  \item[Grande entreprise avec contraintes on-premise fortes]
        Jenkins~: maturité, plugins métier, intégration aux outils existants
        (Active Directory, LDAP, JIRA, Bitbucket).
  \item[Projet cloud-native polyglotte multi-VCS]
        CircleCI~: orbs réutilisables, support GitHub + GitLab + Bitbucket,
        performances élevées.
  \item[Projet OSS hérité encore sur Travis]
        Migrer vers GitHub Actions à la prochaine occasion~; ne plus démarrer
        de nouveau projet sur Travis.
\end{description}

% ══════════════════════════════════════════════════════════════════════════════
\chapter*{Conclusion}
\addcontentsline{toc}{chapter}{Conclusion}

Ce projet a permis de mettre en pratique de manière concrète les concepts de
l'intégration et du déploiement continus sur une application Jakarta EE
réelle. Trois leçons principales peuvent être retenues.

\textbf{Première leçon~-- la testabilité conditionne la CI/CD.} Extraire la
logique métier dans des classes utilitaires pures, prévoir des constructeurs
d'injection dans les ressources REST, ne pas initialiser les connexions à la
base de données dans des blocs statiques sans stratégie de fallback~-- toutes
ces pratiques, inspirées de l'architecture hexagonale, conditionnent
directement la capacité à tester sans infrastructure et donc à intégrer le
code dans un pipeline rapide.

\textbf{Deuxième leçon~-- la portabilité conceptuelle des pipelines.} La même
logique fonctionnelle (lint $\to$ test $\to$ qualité $\to$ package $\to$
docker $\to$ sécurité $\to$ notify) s'exprime sur les cinq moteurs étudiés
avec des syntaxes différentes mais des concepts identiques. Maîtriser un
moteur facilite donc l'adoption des autres. Le choix entre eux relève
davantage de contraintes contextuelles (hébergement Git, politique de sécurité,
budget) que de différences techniques majeures.

\textbf{Troisième leçon~-- la pipeline n'est pas une fin en soi.} Un pipeline
bien configuré mais dont les alertes Slack sont ignorées, dont les rapports
Sonar ne sont jamais consultés, dont les scans Trivy ne déclenchent jamais de
mise à jour, n'offre qu'une illusion de rigueur. La CI/CD est un outil au
service d'une culture d'amélioration continue~; sans cette culture, l'outil
seul ne produit aucune valeur.

Pour ce projet spécifique, \textbf{GitHub Actions} s'est révélé l'outil le
plus adapté~: la configuration en YAML est naturellement lisible, le débogage
en ligne est facilité par les logs détaillés, et l'intégration avec Docker
Hub, GHCR, SonarCloud et Slack ne nécessite aucun plugin supplémentaire. Les
quatre autres pipelines (GitLab, Jenkins, CircleCI, Travis) sont fonctionnels
et démontrent que la même logique se transpose avec un effort modéré.

% ── Bibliographie ─────────────────────────────────────────────────────────────
\chapter*{Bibliographie}
\addcontentsline{toc}{chapter}{Bibliographie}

\begin{enumerate}[leftmargin=*, label={[\arabic*]}]
  \item Fowler, M. (2006). \textit{Continuous Integration}. martinfowler.com.
        \url{https://martinfowler.com/articles/continuousIntegration.html}
  \item Humble, J. \& Farley, D. (2010). \textit{Continuous Delivery~: Reliable
        Software Releases through Build, Test, and Deployment Automation}.
        Addison-Wesley Professional.
  \item GitHub, Inc. (2024). \textit{GitHub Actions Documentation}.
        \url{https://docs.github.com/en/actions}
  \item GitLab B.V. (2024). \textit{GitLab CI/CD Documentation}.
        \url{https://docs.gitlab.com/ee/ci/}
  \item Jenkins Community. (2024). \textit{Jenkins User Documentation}.
        \url{https://www.jenkins.io/doc/}
  \item CircleCI. (2024). \textit{CircleCI Documentation}.
        \url{https://circleci.com/docs/}
  \item Travis CI. (2024). \textit{Travis CI Documentation}.
        \url{https://docs.travis-ci.com/}
  \item SonarSource. (2024). \textit{SonarCloud Documentation}.
        \url{https://docs.sonarcloud.io/}
  \item SonarSource. (2024). \textit{SonarQube Documentation}.
        \url{https://docs.sonarqube.org/}
  \item Aqua Security. (2024). \textit{Trivy Documentation}.
        \url{https://aquasecurity.github.io/trivy/}
  \item Docker, Inc. (2024). \textit{Docker Documentation}.
        \url{https://docs.docker.com/}
  \item Docker, Inc. (2024). \textit{Multi-stage builds}.
        \url{https://docs.docker.com/build/building/multi-stage/}
  \item Oracle. (2024). \textit{Oracle Database Express Edition Container}.
        \url{https://github.com/gvenzl/oci-oracle-xe}
  \item Hibernate ORM 6. (2024). \textit{User Guide}.
        \url{https://docs.jboss.org/hibernate/orm/6.4/userguide/html_single/}
  \item Checkstyle Project. (2024). \textit{Checkstyle Documentation}.
        \url{https://checkstyle.org/}
  \item SpotBugs Project. (2024). \textit{SpotBugs Documentation}.
        \url{https://spotbugs.github.io/}
\end{enumerate}

\end{document}
